\begin{receta}{Tacos de Entraña y Piña}{25 min}
    
    \begin{minipage}[t]{0.35\textwidth}
        \section*{Ingredientes}
        \begin{itemize}[leftmargin=*]
            \item Entraña de ternera
            \item Cebolla dulce
            \item Piña natural
            \item Cilantro fresco
            \item Limas
            \item Tortillas de maíz
            \item Sal y aceite
            \item Guacamole (opcional)
        \end{itemize}
    \end{minipage}
    \hfill
    \begin{minipage}[t]{0.6\textwidth}
        \section*{Preparación}
        \begin{enumerate}
            \item \textbf{La carne:} Cortar la entraña en tiras o tacos pequeños.
            \item \textbf{El sofrito:} Cocinar la carne en la sartén a fuego fuerte hasta que esté dorada. Reservar.
            \item \textbf{El acompañamiento:} Picar muy finamente la cebolla dulce, la piña y el cilantro. Mezclar en un bol con zumo de lima y una pizca de sal.
            \item \textbf{Montaje:} Calentar las tortillas en una sartén limpia.
            \item \textbf{Servir:} Rellenar las tortillas con la carne y la mezcla fresca de piña por encima.
        \end{enumerate}
    \end{minipage}

    \vspace{1em}
    \begin{tcolorbox}[colback=orange!5, colframe=orange!50, title=El rincón del Chef]
        \begin{minipage}{0.25\textwidth}
            \centering
            \usebox{\miqrMexicano} % <--- Aquí llamamos al QR ya renderizado
        \end{minipage}
        \begin{minipage}{0.7\textwidth}
            \textbf{\faMusic\ Música:} Escanea para cocinar con ritmo. \\
            \textbf{\faLightbulb\ Nota:} \textit{Añade una base de guacamole sobre la tortilla caliente antes de poner la carne para que queden más jugosos.}
        \end{minipage}
    \end{tcolorbox}
\end{receta}